% ---------------------------------------------------------------
% Preamble — IEEEtran + XeLaTeX/LuaLaTeX (TeX Live 2022+ / MiKTeX / Overleaf)
% ---------------------------------------------------------------

% Ngôn ngữ và phông chữ
\usepackage{fontspec}
\usepackage{babel}
\babelprovide[main,import]{vietnamese}
% TeX Gyre Termes (giống Times, phông chuẩn của IEEE) hỗ trợ đầy đủ tiếng Việt.
% Gọi theo tên file để chạy được cả trên TeX Live, Overleaf và Tectonic.
\setmainfont{texgyretermes}[Extension=.otf,
  UprightFont=*-regular, BoldFont=*-bold,
  ItalicFont=*-italic, BoldItalicFont=*-bolditalic]
\setsansfont{texgyreheros}[Extension=.otf,
  UprightFont=*-regular, BoldFont=*-bold,
  ItalicFont=*-italic, BoldItalicFont=*-bolditalic]
\setmonofont{DejaVuSansMono}[Extension=.ttf, Scale=0.85,
  BoldFont=*-Bold, ItalicFont=*-Oblique, BoldItalicFont=*-BoldOblique]

% Bảng
\usepackage{booktabs}
\usepackage{tabularx}
\usepackage{multirow}
\usepackage{array}
\usepackage[table]{xcolor}
\newcolumntype{L}[1]{>{\raggedright\arraybackslash}p{#1}}
\newcolumntype{R}{>{\raggedleft\arraybackslash}X}
\newcolumntype{Y}{>{\raggedright\arraybackslash}X}
\renewcommand{\arraystretch}{1.2}

% Hình, sơ đồ, biểu đồ
\usepackage{graphicx}
\usepackage{tikz}
\usetikzlibrary{positioning,arrows.meta,fit,shapes.geometric,shapes.multipart,calc,backgrounds}
\usepackage{pgfplots}
\pgfplotsset{compat=1.18}

% Mã nguồn
\usepackage{listings}
\lstset{
  basicstyle=\ttfamily\scriptsize,
  breaklines=true,
  frame=tb,
  framerule=0.4pt,
  rulecolor=\color{black!40},
  backgroundcolor=\color{black!3},
  numbers=none,
  keywordstyle=\color{blue!60!black}\bfseries,
  commentstyle=\color{green!40!black}\itshape,
  stringstyle=\color{red!50!black},
  columns=fullflexible,
  keepspaces=true,
  showstringspaces=false,
  aboveskip=6pt,
  belowskip=4pt,
  captionpos=t,
  abovecaptionskip=0pt,
  belowcaptionskip=6pt
}
\renewcommand{\lstlistingname}{Đoạn mã}
% IEEEtran in chú thích với khoảng âm nên chú thích của listings bị đè lên khung.
% Dùng chú thích riêng: \codecaption{nội dung}{nhãn} đặt ngay trước lstlisting.
\newcounter{codelisting}
\newcommand{\codecaption}[2]{%
  \par\smallskip\refstepcounter{codelisting}\label{#2}%
  {\centering\footnotesize \lstlistingname~\thecodelisting. #1\par}\nobreak\vspace{2pt}\nobreak}

% Danh sách, liên kết, đường dẫn
\usepackage{enumitem}
\setlist{itemsep=1pt,topsep=3pt,leftmargin=*}
\usepackage{url}
\usepackage[hidelinks]{hyperref}

% Giảm giãn chữ ở cột hẹp khi dòng chứa tên mã/đường dẫn dài
\setlength{\emergencystretch}{2em}
\pgfplotsset{/pgf/number format/1000 sep={.}, /pgf/number format/use comma}

% Danh sách căn trái cho các mục chứa nhiều tên mã
\newenvironment{codeitems}{\begin{itemize}\raggedright}{\end{itemize}}

% Lệnh tiện ích
\newcommand{\code}[1]{\texttt{#1}}
\newcommand{\commitmain}{\code{5dd739f}}
\newcommand{\commitmono}{\code{1e64bfc}}

% Màu dùng chung cho sơ đồ
\definecolor{svc}{HTML}{DCEBFA}
\definecolor{svcborder}{HTML}{2F6DB5}
\definecolor{infra}{HTML}{FCEBD5}
\definecolor{infraborder}{HTML}{C07A1B}
\definecolor{obs}{HTML}{E6F2E1}
\definecolor{obsborder}{HTML}{4C8A3A}
\definecolor{mono}{HTML}{8A8F98}
\definecolor{micro}{HTML}{2F6DB5}

\tikzset{
  box/.style={draw=svcborder,fill=svc,rounded corners=2pt,align=center,
              minimum height=9mm,inner sep=4pt,font=\small},
  infrabox/.style={box,draw=infraborder,fill=infra},
  obsbox/.style={box,draw=obsborder,fill=obs},
  db/.style={cylinder,shape border rotate=90,aspect=0.18,draw=infraborder,
             fill=infra,align=center,minimum height=12mm,inner sep=3pt,font=\small},
  actor/.style={draw=black!60,fill=black!5,rounded corners=6pt,align=center,
                minimum height=8mm,inner sep=4pt,font=\small},
  arr/.style={-{Stealth[length=2.2mm]},thick,draw=black!70},
  dasharr/.style={arr,dashed},
  lbl/.style={font=\scriptsize,fill=white,inner sep=1.5pt}
}
